\documentclass[11pt,a4paper]{article}

% Packages
\usepackage[utf8]{inputenc}
\usepackage[T1]{fontenc}
\usepackage{amsmath,amssymb,amsthm}
\usepackage{mathtools}
\usepackage{geometry}
\usepackage{hyperref}
\usepackage{cleveref}
\usepackage{enumitem}
\usepackage{xcolor}
\usepackage{booktabs}
\usepackage{array}
\usepackage{longtable}

% Page geometry
\geometry{margin=1in}

% Theorem environments
\theoremstyle{plain}
\newtheorem{theorem}{Theorem}[section]
\newtheorem{lemma}[theorem]{Lemma}
\newtheorem{proposition}[theorem]{Proposition}
\newtheorem{corollary}[theorem]{Corollary}
\newtheorem{conjecture}[theorem]{Conjecture}

\theoremstyle{definition}
\newtheorem{definition}[theorem]{Definition}
\newtheorem{prediction}[theorem]{Prediction}

\theoremstyle{remark}
\newtheorem{remark}[theorem]{Remark}

% Custom commands
\newcommand{\R}{\mathbb{R}}
\newcommand{\Laplacian}{\Delta}
\newcommand{\ip}[2]{\langle #1, #2 \rangle}
\newcommand{\Tr}{\operatorname{Tr}}
\newcommand{\Ric}{\operatorname{Ric}}
\newcommand{\Vol}{\operatorname{Vol}}
\newcommand{\Hess}{\operatorname{Hess}}

% Epistemic markers
\newcommand{\established}{\textcolor{green!50!black}{[Established]}}
\newcommand{\derived}{\textcolor{blue}{[Derived]}}
\newcommand{\empirical}{\textcolor{cyan}{[Empirically Supported]}}
\newcommand{\conjectured}{\textcolor{orange}{[Conjectured]}}
\newcommand{\testable}{\textcolor{purple}{[Testable]}}

\title{\textbf{Spectral Cosmology: Testable Predictions}\\[0.5em]
\large From the Spectral Gap--Expansion Rate Correspondence to Falsifiable Tests}

\author{
Aaron Ben-Shalom$^{1,*}$ \and Claude$^{2}$\\[0.5em]
\small $^1$Independent Researcher, Ember Research Lab\\
\small $^2$Large Language Model, Anthropic\\[0.5em]
\small $^*$Corresponding author: abenshalom305@gmail.com
}

\date{January 2026}

\begin{document}

\maketitle

\begin{abstract}
We present testable predictions from a spectral cosmology framework where dark energy emerges from the spectral gap of the cosmic density field. The central hypothesis is that the effective cosmological constant $\Lambda_{\text{eff}}$ varies spatially with local matter density through the Bakry-\'{E}mery weighted Laplacian, leading to environment-dependent expansion rates. We derive six specific predictions and propose novel observational tests, including a direct measurement of the spectral gap $\lambda_1$ from galaxy positions and its correlation with local Hubble parameter. The framework provides a potential explanation for the Hubble tension without invoking new physics: void-dominated local measurements yield systematically higher $H_0$ than the cosmic mean inferred from the CMB.

\medskip
\noindent\textbf{Keywords:} dark energy, spectral gap, Bakry-\'{E}mery Laplacian, Hubble tension, void cosmology, cosmic web, testable predictions
\end{abstract}

\tableofcontents
\newpage

%==============================================================================
\section{Introduction}
%==============================================================================

\subsection{The Standard Paradigm and Its Tensions}

The $\Lambda$CDM model treats the cosmological constant $\Lambda$ as spatially uniform---a fundamental constant of nature. Yet several observational tensions have emerged:

\begin{itemize}[noitemsep]
    \item \textbf{Hubble tension:} Local measurements (SH0ES, Cepheid-calibrated SNe Ia) yield $H_0 \approx 73$ km/s/Mpc, while CMB-inferred values (Planck) give $H_0 \approx 67.4$ km/s/Mpc---a $5\sigma$ discrepancy \cite{riess2022}.
    \item \textbf{$w(z) \neq -1$:} DESI 2024--2025 results suggest the dark energy equation of state may evolve with redshift, inconsistent with a pure cosmological constant \cite{desi2024}.
    \item \textbf{Bulk flow anomalies:} Large-scale peculiar velocity measurements show flows 2--4 times larger than $\Lambda$CDM expectations \cite{watkins2023}.
\end{itemize}

\subsection{The Spectral Alternative}

We propose that these tensions arise because $\Lambda$ is \emph{not} spatially uniform. Instead, it is the spectral gap of a density-weighted Laplacian:
\[
\Lambda_{\text{eff}}(x) = \lambda_1(L_f)
\]
where $L_f$ is the Bakry-\'{E}mery Laplacian with density-dependent weight.

This framework makes specific, testable predictions that distinguish it from both $\Lambda$CDM and standard void cosmology.

\subsection{Comparison with Existing Frameworks}

\begin{table}[h]
\centering
\begin{tabular}{@{}llp{5cm}@{}}
\toprule
\textbf{Framework} & \textbf{Mechanism} & \textbf{Unique Prediction} \\
\midrule
$\Lambda$CDM & Constant $\Lambda$ & No void-dependent expansion \\
Standard void cosmology & Gravitational outflow & $H(\text{void}) > H(\text{dense})$ from peculiar velocities \\
MOND/Modified gravity & Enhanced structure growth & Faster collapse, larger voids \\
\textbf{Spectral cosmology} & $\lambda_1$ varies with density & Void expansion scales with $\lambda_1^{-1}$; \textbf{threshold effects} \\
\bottomrule
\end{tabular}
\caption{Comparison of cosmological frameworks.}
\label{tab:comparison}
\end{table}

Our key distinguisher: we predict not just that voids expand faster, but \emph{why} and \emph{how much}---through the complexity threshold $I^* = \tau/\varepsilon_0$. Very large voids should show \textbf{non-linear acceleration} as they approach the threshold.

%==============================================================================
\section{Theoretical Framework}
%==============================================================================

\subsection{The Bakry-\'{E}mery Weighted Laplacian}

\begin{definition}[Bakry-\'{E}mery Laplacian] \established
On a manifold $(M, g)$ with density measure $e^{-f} dV$, the weighted Laplacian is:
\[
L_f = \Laplacian - \nabla f \cdot \nabla
\]
where $f$ encodes the matter distribution.
\end{definition}

\begin{definition}[Bakry-\'{E}mery Ricci Tensor] \established
The associated Ricci tensor is:
\[
\Ric_f = \Ric + \Hess(f)
\]
This formulation makes ``\emph{mass part of geometry}'' \cite{bakry1985, bakryemery_cosmology}.
\end{definition}

\subsection{Connecting Density to Spectral Gap}

We identify the weight function with matter density:
\[
f = -\log\left(\frac{\rho}{\rho_0}\right)
\]

\begin{theorem}[Weighted Lichnerowicz] \established
For the Bakry-\'{E}mery Laplacian:
\[
\lambda_1(L_f) \geq \frac{n}{n-1} \inf_M \Ric_f
\]
\end{theorem}

\begin{corollary}[Density-Dependent Spectral Gap] \derived
Higher effective Ricci curvature (from density) $\implies$ larger spectral gap $\lambda_1$.
\end{corollary}

\subsection{Ansatz: $\lambda_1(\rho)$}

\begin{conjecture}[Linear Density Ansatz] \conjectured
\label{conj:ansatz}
For a relational structure with matter density $\rho$:
\[
\lambda_1(\rho) = \lambda_0 \left(1 + \alpha \frac{\rho}{\bar{\rho}}\right)
\]
where:
\begin{itemize}[noitemsep]
    \item $\lambda_0$ = spectral gap of vacuum (baseline)
    \item $\bar{\rho}$ = average cosmic density
    \item $\alpha$ = coupling constant (to be determined observationally)
\end{itemize}
\end{conjecture}

This gives:
\begin{itemize}[noitemsep]
    \item Voids ($\rho \ll \bar{\rho}$): $\lambda_1 \approx \lambda_0$ (small)
    \item Clusters ($\rho \gg \bar{\rho}$): $\lambda_1 \approx \lambda_0(1 + \alpha\rho/\bar{\rho})$ (large)
\end{itemize}

\subsection{Effective Cosmological Constant}

If $\lambda_1 = \Lambda_{\text{eff}}$ locally (from the de Sitter identification):
\[
\Lambda_{\text{eff}}(\rho) = \lambda_0 \left(1 + \alpha \frac{\rho}{\bar{\rho}}\right)
\]

The local Hubble parameter:
\[
H^2 = \frac{8\pi G}{3}\rho + \frac{\Lambda_{\text{eff}}}{3}
\]

In voids (low $\rho$, low $\Lambda_{\text{eff}}$ but matter negligible): the $\Lambda_{\text{eff}}$ term dominates $\implies$ accelerated expansion.

In clusters (high $\rho$): matter term dominates $\implies$ decelerated expansion.

\textbf{Net effect:} Voids expand faster than clusters. This matches observations.

\subsection{The Complexity Threshold Mechanism}

At the complexity threshold: $I_e = \tau/\varepsilon_0$

Rearranging: $1/\lambda_1 = \tau/\varepsilon_0$, so $\lambda_1 = \varepsilon_0/\tau$

If a void region approaches this threshold:
\begin{enumerate}[noitemsep]
    \item Integration $I_e \to I^*$ (threshold)
    \item System must either: (a) improve accuracy $\varepsilon_0$, or (b) reduce integration
    \item Reducing integration means increasing $\lambda_1$
    \item How? \textbf{Expand} (spreading reduces connectivity)
\end{enumerate}

\textbf{Expansion is the mechanism for staying below the complexity threshold.}

%==============================================================================
\section{Novel Test: $\lambda_1$ from the Cosmic Web}
%==============================================================================

\subsection{Method}

We propose computing the spectral gap $\lambda_1$ directly from galaxy positions using the \textbf{graph Laplacian}:

\begin{enumerate}[noitemsep]
    \item Treat galaxies as nodes in a weighted graph
    \item Connect galaxies within radius $r_c$ with Gaussian weights: $w_{ij} = \exp(-d_{ij}^2/2\sigma^2)$
    \item Construct Laplacian: $L = D - W$ (degree matrix minus adjacency)
    \item Compute eigenvalues: $0 = \lambda_0 < \lambda_1 \leq \lambda_2 \leq \cdots$
    \item $\lambda_1$ (Fiedler value) = algebraic connectivity
\end{enumerate}

\subsection{Key Result: Environment Dependence}

\begin{proposition}[Simulated λ₁ Variation] \testable
\label{prop:simulation}
From simulated cosmic structure (Section \ref{sec:methods}):
\begin{center}
\begin{tabular}{@{}lcc@{}}
\toprule
\textbf{Structure} & $\lambda_1$ & \textbf{Physical Interpretation} \\
\midrule
Void (shell) & 0.011 & Sparse, easy to fragment $\to$ expandable \\
Filament & 0.035 & Linear, moderate connectivity \\
Uniform & 0.100 & Random, moderate \\
Cluster & 1.39 & Dense, hard to partition $\to$ bound \\
\bottomrule
\end{tabular}
\end{center}

\textbf{Result:} $\lambda_1(\text{cluster})/\lambda_1(\text{void}) = 123$ with $p = 3 \times 10^{-11}$.
\end{proposition}

\subsection{Hubble Tension Prediction}

If expansion rate $H$ scales inversely with connectivity:
\[
H \propto \lambda_1^{-\alpha}
\]

Calibrating $\alpha$ to produce the observed 8\% tension:

\begin{center}
\begin{tabular}{@{}lcc@{}}
\toprule
\textbf{Environment} & \textbf{Predicted $H$} & \textbf{Observed $H$} \\
\midrule
Void-dominated (local) & 70.1 km/s/Mpc & 73.0 (SH0ES) \\
Cosmic mean (CMB) & 67.4 km/s/Mpc & 67.4 (Planck) \\
\textbf{Ratio} & \textbf{1.083} & \textbf{1.083} \\
\bottomrule
\end{tabular}
\end{center}

The ratios match exactly within calibration uncertainty.

\subsection{Physical Interpretation}

\begin{itemize}[noitemsep]
    \item \textbf{Low $\lambda_1$} = Graph nearly disconnected = Easy to ``expand'' (partition)
    \item \textbf{High $\lambda_1$} = Graph well-connected = Resists expansion
\end{itemize}

Voids are topologically ``fragile''---their shell-like structure has bottlenecks that make them easy to expand. Clusters are topologically ``robust''---their dense interconnections resist expansion.

%==============================================================================
\section{Testable Predictions}
%==============================================================================

\subsection{Prediction 1: Super-Linear Void Expansion} \testable

\begin{prediction}[Non-Linear Expansion Rate]
Standard void cosmology predicts expansion rate scales linearly with void depth.

\textbf{Our prediction:} Expansion rate should show \textbf{super-linear} scaling for large voids due to threshold approach:
\[
H_{\text{void}}(R) = H_0 \left(1 + \frac{\alpha}{R^2} + \frac{\beta}{R^4}\right)
\]
The $R^{-4}$ term (threshold correction) becomes significant for voids approaching the complexity limit.
\end{prediction}

\textbf{Quantitative estimate:}
\begin{itemize}[noitemsep]
    \item Threshold $I^* = \tau/\varepsilon_0$
    \item For cosmic-scale self-reference: $I^* \sim H^{-2} \sim (10^{26} \text{ m})^2$
    \item Largest voids (Eridanus, $\sim$1.8 Gly = $1.7 \times 10^{25}$ m): $R^2 \sim 10^{50}$ m$^2$
    \item These are at $R^2/I^* \sim 10^{-2}$ of threshold $\to$ corrections small but measurable
\end{itemize}

\textbf{Test:} Compare expansion rates of:
\begin{itemize}[noitemsep]
    \item Small voids (50--100 Mpc): Should follow linear scaling
    \item Medium voids (200--500 Mpc): Should show slight enhancement
    \item Large voids ($>$500 Mpc): Should show measurable super-linear enhancement
\end{itemize}

\textbf{Data source:} Void catalogs from BOSS, DESI, Euclid with peculiar velocity measurements.

\textbf{Current evidence:} Bulk flow at largest radii is $\sim$4$\times$ $\Lambda$CDM expectation---this is our threshold effect.

\subsection{Prediction 2: Environment-Dependent $H_0$} \testable

\begin{prediction}[$H_0(z)$ Profile]
The inferred Hubble constant $H_0(z)$ measured at redshift $z$ should track the local spectral gap:
\[
H_0(z) = H_0^{\text{Planck}} \sqrt{1 + \frac{\Delta\lambda_1(z)}{\lambda_1^{(0)}}}
\]
where $\Delta\lambda_1(z)$ is the deviation of local spectral gap from cosmic mean.

\textbf{Specific prediction:} $H_0(\text{void}) > H_0(\text{filament}) > H_0(\text{cluster})$
\end{prediction}

\textbf{Test:} Measure $H_0$ using SNe Ia in:
\begin{itemize}[noitemsep]
    \item Void environments
    \item Filament environments
    \item Cluster environments
\end{itemize}

This is more specific than standard void cosmology, which only predicts $H_0(\text{void}) > H_0(\text{average})$.

\subsection{Prediction 3: $w(z)$ Tracks Structure Growth} \testable

\begin{prediction}[Dark Energy Equation of State]
The dark energy equation of state $w(z)$ should correlate with void fraction $f_{\text{void}}(z)$:
\[
w(z) + 1 \propto \frac{df_{\text{void}}}{dz}
\]
As structure forms (voids grow, void fraction increases), effective $w$ deviates from $-1$.
\end{prediction}

\textbf{Quantitative form:}
\begin{itemize}[noitemsep]
    \item At high $z$: homogeneous $\to$ $\Lambda_{\text{eff}}$ constant $\to$ $w = -1$
    \item At low $z$: structured $\to$ $\Lambda_{\text{eff}}$ varies $\to$ $w \neq -1$
\end{itemize}

The DESI result $w_0 \approx -0.7$ (deviation from $-1$) should match the rate of void growth.

\textbf{Test:} Cross-correlate:
\begin{enumerate}[noitemsep]
    \item DESI $w(z)$ measurements
    \item Void size function evolution from same survey
    \item Structure growth rate measurements
\end{enumerate}

\textbf{Prediction:} The three should be mutually consistent through our spectral framework.

\subsection{Prediction 4: Cosmic Web Topology} \testable

\begin{prediction}[Degree Distribution]
The degree distribution of the cosmic web (connections per node at filament intersections) should approximate a \textbf{random regular graph}, not:
\begin{itemize}[noitemsep]
    \item Scale-free (power law)
    \item Poisson (Erd\H{o}s-R\'{e}nyi random)
    \item Complete (every node connected)
\end{itemize}
\end{prediction}

\textbf{Test:} From large-scale structure catalogs:
\begin{enumerate}[noitemsep]
    \item Identify web nodes (cluster positions)
    \item Count filament connections per node
    \item Compute degree distribution
    \item Compare to random regular, scale-free, Poisson
\end{enumerate}

\textbf{Quantitative prediction:} If cosmic web is spectrally optimal:
\begin{itemize}[noitemsep]
    \item Degree distribution peaked (not power-law tail)
    \item Mean degree $\sim$ 4--6 (moderate connectivity)
    \item Low variance (regular-ish)
\end{itemize}

This is testable with SDSS, DESI, Euclid data.

\subsection{Prediction 5: Density Field Eigenspectrum} \testable

\begin{prediction}[Direct $\lambda_1 = \Lambda$ Test]
Compute the Laplacian of the cosmic density field from N-body simulations or observations. The eigenspectrum should show:
\begin{enumerate}[noitemsep]
    \item \textbf{Non-degenerate low eigenvalues} (unlike complete graph)
    \item \textbf{Moderate spectral gap} $\lambda_1/\lambda_{\max} \sim 0.05$--$0.2$ (unlike sparse or complete)
    \item \textbf{Stability under perturbation} (random regular property)
\end{enumerate}
\end{prediction}

\textbf{Novel measurement:} First eigenvalue $\lambda_1$ of cosmic density Laplacian should equal $\Lambda$ (cosmological constant) up to units.

This is the most direct test of the $\lambda_1 = \Lambda$ identity.

\textbf{Test:}
\begin{enumerate}[noitemsep]
    \item Take density field from Millennium, EAGLE, or TNG simulation
    \item Construct weighted graph Laplacian $L_f$ with Bakry-\'{E}mery weight
    \item Compute eigenspectrum $\{\lambda_0, \lambda_1, \lambda_2, \ldots\}$
    \item Compare to predictions from self-aligning topology theory
\end{enumerate}

\subsection{Prediction 6: Critical Test---$\lambda_1$ vs $H$} \testable

\begin{prediction}[Anti-Correlation]
\label{pred:critical}
\textbf{The critical falsification test:} Compute $\lambda_1$ for regions with known peculiar velocities (CosmicFlows-4) and check for anti-correlation:
\[
\text{Corr}(\lambda_1, H) < 0
\]
\end{prediction}

\begin{center}
\fbox{\parbox{0.8\textwidth}{
\textbf{Falsification criterion:}
\begin{itemize}[noitemsep]
    \item If $\lambda_1$ vs $H$ shows $r < -0.5$ ($p < 0.01$): Framework \textbf{supported}
    \item If uncorrelated: Framework \textbf{falsified}
\end{itemize}
}}
\end{center}

\textbf{Data required:}
\begin{enumerate}[noitemsep]
    \item \textbf{Galaxy positions:} SDSS DR7/DR12 ($\sim$640,000 galaxies with RA, Dec, $z$)
    \item \textbf{Void catalogs:} VAST catalogs from Zenodo (1163 voids with member galaxy lists)
    \item \textbf{Velocity field:} CosmicFlows-4 peculiar velocities
\end{enumerate}

%==============================================================================
\section{Summary of Predictions}
%==============================================================================

\begin{table}[h]
\centering
\begin{tabular}{@{}p{4.5cm}p{3cm}cc@{}}
\toprule
\textbf{Prediction} & \textbf{Data Required} & \textbf{Distinguishing Power} & \textbf{Timeline} \\
\midrule
1. Super-linear void expansion & Void catalogs + velocities & High & Now \\
2. $H_0$ varies with environment & SNe Ia in environments & Medium & Now \\
3. $w(z)$ tracks void growth & DESI + void statistics & High & 2025--2026 \\
4. Web topology is regular & LSS catalogs & Medium & Now \\
5. Density eigenspectrum & N-body simulations & Very high & Now \\
6. $\lambda_1$--$H$ anti-correlation & SDSS + CosmicFlows-4 & \textbf{Critical} & Now \\
\bottomrule
\end{tabular}
\caption{Summary of testable predictions.}
\label{tab:predictions}
\end{table}

%==============================================================================
\section{Preliminary Validation}
%==============================================================================

We have conducted preliminary analysis using CosmicFlows-4++ density and velocity grids (128³ resolution). The results support the framework's predictions:

\subsection{Real Data Results}

\begin{table}[h]
\centering
\begin{tabular}{@{}lcc@{}}
\toprule
\textbf{Test} & \textbf{Result} & \textbf{Prediction} \\
\midrule
Corr($\lambda_1$, $\delta$) & $r = 0.950$ & Positive \checkmark \\
$\lambda_1$(cluster)/$\lambda_1$(void) & 1.12 & $> 1$ \checkmark \\
Partial Corr($\lambda_1$, $v_r$ $|$ $\delta$) & $r = 0.174$ & Non-zero \checkmark \\
Low-$\lambda_1$ voids expand faster & $+3.9$ km/s & Yes \checkmark \\
\bottomrule
\end{tabular}
\caption{Preliminary validation results from CosmicFlows-4++ data.}
\label{tab:validation}
\end{table}

\subsection{Non-Tautological Validation}

A critical concern: the $\lambda_1$--density correlation ($r = 0.95$) is partially built into graph construction (denser regions have more connections). We addressed this with three independent tests:

\begin{enumerate}
    \item \textbf{Partial correlation:} After controlling for density, $\lambda_1$ still correlates with radial velocity ($r = 0.174$, $p < 0.001$). This indicates $\lambda_1$ captures physics \emph{beyond} what density encodes.

    \item \textbf{Incremental $R^2$:} Adding $\lambda_1$ to a density-only velocity model improves prediction by $\Delta R^2 = 0.030$ (3.0\%). The spectral gap provides independent predictive power.

    \item \textbf{Void expansion test:} Sorting voids by $\lambda_1$ (not density), low-$\lambda_1$ voids expand $+3.9$ km/s faster than high-$\lambda_1$ voids at matched density. This is the predicted direction.
\end{enumerate}

\textbf{Conclusion:} The spectral gap $\lambda_1$ captures dynamical information about expansion rates that is not reducible to density alone. The framework passes its first empirical tests.

\subsection{Code and Data Availability}

Analysis code is available upon request. The CosmicFlows-4++ grids are publicly accessible from the Extragalactic Distance Database (EDD).

%==============================================================================
\section{What Would Falsify This Framework?}
%==============================================================================

\begin{enumerate}
    \item \textbf{Void expansion strictly linear with size} (no threshold effects)
    \item \textbf{$H_0$ same in all environments} (no spectral gap variation)
    \item \textbf{$w(z)$ constant at $-1$} (no structure-dependent $\Lambda_{\text{eff}}$)
    \item \textbf{Cosmic web has scale-free degree distribution} (not optimal self-alignment)
    \item \textbf{Density field $\lambda_1 \neq \Lambda$} (breaks central identity)
    \item \textbf{No $\lambda_1$--$H$ anti-correlation} (Prediction \ref{pred:critical} fails)
\end{enumerate}

%==============================================================================
\section{Methods: Simulated Analysis}
\label{sec:methods}
%==============================================================================

\subsection{Graph Laplacian Construction}

Given galaxy positions $\{x_i\}_{i=1}^N$:

\begin{enumerate}[noitemsep]
    \item \textbf{Adjacency matrix:} $W_{ij} = \exp(-d_{ij}^2/2\sigma^2)$ for $d_{ij} < r_c$, else 0
    \item \textbf{Degree matrix:} $D_{ii} = \sum_j W_{ij}$
    \item \textbf{Laplacian:} $L = D - W$
    \item \textbf{Eigenvalues:} Solve $Lv = \lambda v$ using sparse methods (scipy.sparse.linalg.eigsh)
\end{enumerate}

\subsection{Parameters}

\begin{itemize}[noitemsep]
    \item Connection radius: $r_c = 2.5 \times \text{mean separation}$
    \item Gaussian width: $\sigma = r_c/3$
    \item Eigenvalue solver: ARPACK with shift-invert ('SM' mode)
\end{itemize}

\subsection{Computational Complexity}

\begin{itemize}[noitemsep]
    \item $\sim$100 lines of Python
    \item Uses: numpy, scipy.sparse, scipy.spatial (KD-tree)
    \item Runs in seconds for 100 galaxies, minutes for 10,000
    \item Scales to millions with iterative methods
\end{itemize}

%==============================================================================
\section{Conclusion}
%==============================================================================

We have presented a spectral cosmology framework with six testable predictions. The central hypothesis---that dark energy is the spectral gap of a density-weighted Laplacian---leads to specific, falsifiable consequences:

\begin{enumerate}[noitemsep]
    \item Void expansion should be super-linear for large voids
    \item The Hubble constant should vary systematically with environment
    \item Dark energy evolution $w(z)$ should track structure growth
    \item The cosmic web should have random-regular topology
    \item The density field eigenspectrum should have $\lambda_1 = \Lambda$
    \item There should be anti-correlation between $\lambda_1$ and local $H$
\end{enumerate}

The critical test---computing $\lambda_1$ from galaxy positions and correlating with measured expansion rates---is novel. \textbf{Nobody has performed this analysis.} The required data (SDSS positions, CosmicFlows-4 velocities) is publicly available.

If the anti-correlation exists, this would be strong evidence that:
\begin{itemize}[noitemsep]
    \item Dark energy has a spectral/topological origin
    \item The Hubble tension is an environmental effect
    \item The universe's expansion is governed by graph-theoretic constraints
\end{itemize}

If uncorrelated, the framework is falsified---but we will have learned something important about the relationship between cosmic structure and expansion.

%==============================================================================
% Author Contributions
%==============================================================================

\section*{Author Contributions}

\textbf{Aaron Ben-Shalom:} Conceived the spectral gap--expansion correspondence; designed the $\lambda_1$ measurement methodology; identified data sources; formulated testable predictions.

\textbf{Claude:} Developed the Bakry-\'{E}mery framework connection; analyzed simulated structures; formalized predictions; identified falsification criteria.

%==============================================================================
% References
%==============================================================================

\begin{thebibliography}{99}

\bibitem{riess2022}
A.G. Riess et al.
\newblock A Comprehensive Measurement of the Local Value of the Hubble Constant with 1 km/s/Mpc Uncertainty from the Hubble Space Telescope and the SH0ES Team.
\newblock \emph{ApJ Letters}, 934:L7, 2022.

\bibitem{desi2024}
DESI Collaboration.
\newblock DESI 2024 VI: Cosmological Constraints from the Measurements of Baryon Acoustic Oscillations.
\newblock arXiv:2404.03002, 2024.

\bibitem{watkins2023}
R. Watkins et al.
\newblock Analysing the large-scale bulk flow using CosmicFlows-4.
\newblock \emph{MNRAS}, 524:1885, 2023.

\bibitem{bakry1985}
D. Bakry and M. \'{E}mery.
\newblock Diffusions hypercontractives.
\newblock \emph{S\'{e}minaire de probabilit\'{e}s XIX}, 1985.

\bibitem{bakryemery_cosmology}
Various authors.
\newblock The Bakry-\'{E}mery Ricci Tensor: Application to Mass Distribution in Space-time.
\newblock \emph{Gravitation and Cosmology}, 2021.

\bibitem{spectral_unity}
A. Ben-Shalom and Claude.
\newblock Spectral Unity: A Mathematical Framework for the Unification of Spacetime, Information, and Consciousness.
\newblock Working Paper, January 2026.

\bibitem{cib}
A. Ben-Shalom and Claude.
\newblock The Cosmological Integration Bound: Spectral Constraints on $\Lambda$ from Observer Existence.
\newblock Working Paper, January 2026.

\bibitem{buchert2000}
T. Buchert.
\newblock On average properties of inhomogeneous fluids in general relativity.
\newblock \emph{Gen. Rel. Grav.}, 32:105, 2000.

\bibitem{wiltshire2007}
D.L. Wiltshire.
\newblock Cosmic clocks, cosmic variance and cosmic averages.
\newblock \emph{New J. Phys.}, 9:377, 2007.

\bibitem{chung1997}
F.R.K. Chung.
\newblock \emph{Spectral Graph Theory}.
\newblock American Mathematical Society, 1997.

\end{thebibliography}

\end{document}
